\documentclass[../main.tex]{subfiles}
\begin{document}

\chapter{Threads and Concurrency}
\lessoninfo{
  video={P2L2},
  textbook={OSTEP \cite{ostep} ch.~26--32; Silberschatz et al.\ \cite{silberschatz2018} ch.~4--5, 8},
  keywords={thread, multithreading, concurrency, race condition, mutex, condition variable, readers-writers problem, spurious wake-up, deadlock, kernel-level thread, user-level thread, multithreading models, boss-workers pattern, pipeline pattern},
}

A process as described in Lesson~2 has a single execution context, so it
can follow only one sequence of instructions at a time.  This lesson
introduces threads, which give a process several execution contexts that
share one address space, and the mechanisms---mutexes and condition
variables---with which such threads coordinate their access to shared data.
The presentation follows Birrell's tutorial on programming with threads
\cite{birrell1989}; it then turns to deadlocks, to the relation between
user-level and kernel-level threads, and to common patterns for structuring
multithreaded applications.

\section{Processes and threads}
\label{sec:l03-threads}

A process is represented by its address space---the code, data and heap,
reached through the virtual-to-physical mappings kept in the page
table---, by the resources the operating system keeps for it, such as open
files, and by its execution context: the values of the CPU registers,
including the program counter and the stack pointer, and the
stack.\source{P2L2, process vs.\ thread; Birrell \cite{birrell1989}; OSTEP
ch.~26.}  The operating system records all of this in the process control
block.  A process of this kind is single-threaded.

\begin{definition}[Thread]
  A \term{thread}[スレッド] is one execution context within a process.  Every
  thread has its own program counter, stack pointer, register values and
  stack; all threads of a process share its address space (code, data,
  heap) and its other resources, such as open files.
\end{definition}

Running several threads inside one process is called
\term{multithreading}[マルチスレッド], and such a process is
\emph{multithreaded}.  Since the threads share the virtual-to-physical
mappings, all of them can access the same code and data.  At any moment,
however, they may execute different instructions, work on different
portions of the input and have different chains of procedure calls in
progress.  The per-thread state in \cref{fig:l03-process-thread} captures
exactly these differences: the program counter and the registers record
where a thread is and what it is computing, and the private stack holds the
local variables and return addresses of its active calls.

\begin{figure}[htbp]
  \centering
  % Single-threaded vs. multithreaded process: shared and per-thread state.
% Usage: % Single-threaded vs. multithreaded process: shared and per-thread state.
% Usage: % Single-threaded vs. multithreaded process: shared and per-thread state.
% Usage: \input{figures/l03-process-thread}   (width ~116 mm)
\begin{tikzpicture}[x=1mm, y=1mm, font=\footnotesize\sffamily,
  sh/.style={draw, fill=ln-box, minimum height=6mm, inner sep=0pt, anchor=south west},
  pt/.style={draw, fill=white, minimum height=6mm, minimum width=18mm,
             inner sep=0pt, anchor=south west},
  run/.style={-{Stealth[length=2mm]}, very thick, ln-accent},
  frame/.style={draw=ln-rule, rounded corners=2pt},
  ttl/.style={font=\footnotesize\sffamily\bfseries, anchor=south}]
  % ---- single-threaded process
  \draw[frame] (-2,-12) rectangle (44,22);
  \node[ttl] at (21,23) {\iflangja{シングルスレッドプロセス}{single-threaded process}};
  \node[sh, minimum width=14mm] at (0,14)  {\iflangja{コード}{code}};
  \node[sh, minimum width=14mm] at (14,14) {\iflangja{データ}{data}};
  \node[sh, minimum width=14mm] at (28,14) {\iflangja{ファイル}{files}};
  \node[pt] at (12,7) {\iflangja{レジスタ}{registers}};
  \node[pt] at (12,0) {\iflangja{スタック}{stack}};
  \draw[run] (21,-1.5) -- (21,-10);
  \node[anchor=west, text=ln-accent] at (22.5,-6) {\iflangja{スレッド}{thread}};
  % ---- multithreaded process
  \draw[frame] (52,-12) rectangle (116,22);
  \node[ttl] at (84,23) {\iflangja{マルチスレッドプロセス}{multithreaded process}};
  \node[sh, minimum width=20mm] at (54,14) {\iflangja{コード}{code}};
  \node[sh, minimum width=20mm] at (74,14) {\iflangja{データ}{data}};
  \node[sh, minimum width=20mm] at (94,14) {\iflangja{ファイル}{files}};
  \foreach \x in {54,75,96} {
    \node[pt] at (\x,7) {\iflangja{レジスタ}{registers}};
    \node[pt] at (\x,0) {\iflangja{スタック}{stack}};
    \draw[run] (\x+9,-1.5) -- (\x+9,-10);
  }
\end{tikzpicture}
   (width ~116 mm)
\begin{tikzpicture}[x=1mm, y=1mm, font=\footnotesize\sffamily,
  sh/.style={draw, fill=ln-box, minimum height=6mm, inner sep=0pt, anchor=south west},
  pt/.style={draw, fill=white, minimum height=6mm, minimum width=18mm,
             inner sep=0pt, anchor=south west},
  run/.style={-{Stealth[length=2mm]}, very thick, ln-accent},
  frame/.style={draw=ln-rule, rounded corners=2pt},
  ttl/.style={font=\footnotesize\sffamily\bfseries, anchor=south}]
  % ---- single-threaded process
  \draw[frame] (-2,-12) rectangle (44,22);
  \node[ttl] at (21,23) {\iflangja{シングルスレッドプロセス}{single-threaded process}};
  \node[sh, minimum width=14mm] at (0,14)  {\iflangja{コード}{code}};
  \node[sh, minimum width=14mm] at (14,14) {\iflangja{データ}{data}};
  \node[sh, minimum width=14mm] at (28,14) {\iflangja{ファイル}{files}};
  \node[pt] at (12,7) {\iflangja{レジスタ}{registers}};
  \node[pt] at (12,0) {\iflangja{スタック}{stack}};
  \draw[run] (21,-1.5) -- (21,-10);
  \node[anchor=west, text=ln-accent] at (22.5,-6) {\iflangja{スレッド}{thread}};
  % ---- multithreaded process
  \draw[frame] (52,-12) rectangle (116,22);
  \node[ttl] at (84,23) {\iflangja{マルチスレッドプロセス}{multithreaded process}};
  \node[sh, minimum width=20mm] at (54,14) {\iflangja{コード}{code}};
  \node[sh, minimum width=20mm] at (74,14) {\iflangja{データ}{data}};
  \node[sh, minimum width=20mm] at (94,14) {\iflangja{ファイル}{files}};
  \foreach \x in {54,75,96} {
    \node[pt] at (\x,7) {\iflangja{レジスタ}{registers}};
    \node[pt] at (\x,0) {\iflangja{スタック}{stack}};
    \draw[run] (\x+9,-1.5) -- (\x+9,-10);
  }
\end{tikzpicture}
   (width ~116 mm)
\begin{tikzpicture}[x=1mm, y=1mm, font=\footnotesize\sffamily,
  sh/.style={draw, fill=ln-box, minimum height=6mm, inner sep=0pt, anchor=south west},
  pt/.style={draw, fill=white, minimum height=6mm, minimum width=18mm,
             inner sep=0pt, anchor=south west},
  run/.style={-{Stealth[length=2mm]}, very thick, ln-accent},
  frame/.style={draw=ln-rule, rounded corners=2pt},
  ttl/.style={font=\footnotesize\sffamily\bfseries, anchor=south}]
  % ---- single-threaded process
  \draw[frame] (-2,-12) rectangle (44,22);
  \node[ttl] at (21,23) {\iflangja{シングルスレッドプロセス}{single-threaded process}};
  \node[sh, minimum width=14mm] at (0,14)  {\iflangja{コード}{code}};
  \node[sh, minimum width=14mm] at (14,14) {\iflangja{データ}{data}};
  \node[sh, minimum width=14mm] at (28,14) {\iflangja{ファイル}{files}};
  \node[pt] at (12,7) {\iflangja{レジスタ}{registers}};
  \node[pt] at (12,0) {\iflangja{スタック}{stack}};
  \draw[run] (21,-1.5) -- (21,-10);
  \node[anchor=west, text=ln-accent] at (22.5,-6) {\iflangja{スレッド}{thread}};
  % ---- multithreaded process
  \draw[frame] (52,-12) rectangle (116,22);
  \node[ttl] at (84,23) {\iflangja{マルチスレッドプロセス}{multithreaded process}};
  \node[sh, minimum width=20mm] at (54,14) {\iflangja{コード}{code}};
  \node[sh, minimum width=20mm] at (74,14) {\iflangja{データ}{data}};
  \node[sh, minimum width=20mm] at (94,14) {\iflangja{ファイル}{files}};
  \foreach \x in {54,75,96} {
    \node[pt] at (\x,7) {\iflangja{レジスタ}{registers}};
    \node[pt] at (\x,0) {\iflangja{スタック}{stack}};
    \draw[run] (\x+9,-1.5) -- (\x+9,-10);
  }
\end{tikzpicture}

  \caption{A single-threaded and a multithreaded process.  The shaded state
    (code, data and open files) is shared by all threads; every thread has
    its own registers---including PC and SP---and its own stack.}
  \label{fig:l03-process-thread}
\end{figure}

The operating system's representation of a multithreaded process is
correspondingly richer than a PCB for a single-threaded one.  It holds the
information shared by all threads once, and a separate data structure for
each thread that records the thread's execution context.

\section{Why threads are useful}
\label{sec:l03-benefits}

On a machine with several CPUs, threads can speed up an application in two
ways.\source{P2L2, why are threads useful, benefits of multithreading on a
single CPU, benefits to applications and OS code; OSTEP ch.~26.}

\begin{description}
  \item[Same code, different data] All threads execute the same code, each
    on a different portion of the input.  In a matrix multiplication, for
    example, each thread can compute a different subset of the rows of the
    result, and the threads run in parallel on different CPUs.  Dividing the
    work of one task in this way is called
    \term{parallelization}[並列化].
  \item[Different code] Different threads execute different parts of the
    program: one thread may process input while another renders the
    display, or threads may be dedicated to requests of different types.
    Dedicating threads to particular tasks is called
    \term{specialization}[特化].
\end{description}

Specialization has two further benefits.  The threads can be managed
differently, for example by giving higher priority to the threads that
handle more important tasks or customers.  And since performance depends on
how much of a thread's state is present in the processor cache, a thread
that repeatedly executes a small portion of the code works with a hot cache
(Lesson~2) more often than a thread that executes everything.

\subsection{Threads versus processes}

The same parallelism could be obtained by running several single-threaded
processes.  A multithreaded implementation is nevertheless more efficient in
three respects:
\begin{itemize}
  \item \textbf{Memory.}  The threads share one address space, whereas every
    process needs its own address space and its own copy of the state.  The
    multithreaded application has lower memory requirements and is
    therefore less likely to need swapping.
  \item \textbf{Context switching.}  Switching between threads of the same
    process does not switch the address space: the virtual-to-physical
    mappings stay in place, so the costly part of a process context switch
    is avoided.
  \item \textbf{Communication.}  Threads exchange data through shared
    variables.  Processes need inter-process communication (IPC): message
    passing enters the kernel for every exchange, and shared memory must
    first be set up by the operating system, so communication between
    processes is more costly.
\end{itemize}

\begin{example}
  Suppose an application holds \SI{400}{\mega\byte} of data that each
  process would have to keep separately (code pages, which the operating
  system can share between processes, are ignored here) and needs about
  \SI{1}{\mega\byte} of stack and other per-context state.\margin{On Linux a
  thread's stack is \SI{8}{\mebi\byte} of address space by default
  (\texttt{RLIMIT\_STACK}); only the pages it touches occupy memory, so
  \num{10000} threads reserve \SI{80}{\gibi\byte} and may still need no more
  than tens of megabytes.}  Run as four
  processes, each with its own copy of the data, it occupies about
  $4\times 401 = \SI{1604}{\mega\byte}$; run as four threads of one
  process it occupies about $400 + 4\times 1 = \SI{404}{\mega\byte}$.
\end{example}

\subsection{Threads on a single CPU}

Threads are useful even when there are more threads than CPUs.  When a
thread issues a request that blocks, such as a disk read, it has nothing to
do until the request completes, and the CPU would sit idle for that time
$t_{\mathrm{idle}}$.  The CPU can instead switch to another thread and
switch back when the request completes (\cref{fig:l03-idle-time}).  The
other thread then gains $t_{\mathrm{idle}} - t_{\mathrm{ctx}}$ of CPU time,
while the blocked thread resumes $t_{\mathrm{ctx}}$ later than it would
have; the two switches together cost $2\,t_{\mathrm{ctx}}$.  Switching pays
off when the time gained exceeds the delay, that is, when
\begin{equation}
  t_{\mathrm{idle}} \;>\; 2\, t_{\mathrm{ctx}} .
  \label{eq:l03-idle}
\end{equation}
Because a context switch between threads avoids the address-space switch,
its $t_{\mathrm{ctx}}$ is smaller than that of a switch between processes,
and \cref{eq:l03-idle} holds for shorter idle periods.  Multithreading
therefore lets a program hide idle time even on one CPU.\margin{A disk seek
(some \SI{10}{\milli\second}) or an SSD read (some \SI{100}{\micro\second})
repays two switches many times over; a wait of the order of a memory access
(under \SI{100}{\nano\second}) never does, and such short waits are spun on
instead (Lesson~10).}

\begin{figure}[htbp]
  \centering
  % Hiding I/O idle time on a single CPU by switching to another thread.
% Usage: % Hiding I/O idle time on a single CPU by switching to another thread.
% Usage: % Hiding I/O idle time on a single CPU by switching to another thread.
% Usage: \input{figures/l03-idle-time}   (width ~112 mm)
\begin{tikzpicture}[x=1mm, y=1mm, font=\footnotesize\sffamily,
  ev/.style={-{Stealth[length=1.6mm]}},
  dim/.style={{Bar[width=2mm]}-{Bar[width=2mm]}, thin},
  cs/.style={draw, fill=ln-accent!55}]
  \node[anchor=east] at (-1,3) {CPU};
  \draw[fill=ln-box] (0,0) rectangle (28,6);   \node at (14,3) {T1};
  \draw[cs]          (28,0) rectangle (33,6);
  \draw[fill=white]  (33,0) rectangle (83,6);  \node at (58,3) {T2};
  \draw[cs]          (83,0) rectangle (88,6);
  \draw[fill=ln-box] (88,0) rectangle (110,6); \node at (99,3) {T1};
  % events
  \draw[ev] (28,13) node[above] {\iflangja{T1 が I/O でブロック}{T1 blocks on I/O}} -- (28,7);
  \draw[ev] (83,13) node[above] {\iflangja{I/O 完了}{I/O completes}} -- (83,7);
  % dimensions
  \draw[dim] (28,-3) -- (33,-3);
  \draw[dim] (83,-3) -- (88,-3);
  \node[below] at (30.5,-3.5) {$t_{\mathrm{ctx}}$};
  \node[below] at (85.5,-3.5) {$t_{\mathrm{ctx}}$};
  \draw[dim] (28,-11) -- (83,-11);
  \node[below] at (55.5,-11.5) {$t_{\mathrm{idle}}$};
\end{tikzpicture}
   (width ~112 mm)
\begin{tikzpicture}[x=1mm, y=1mm, font=\footnotesize\sffamily,
  ev/.style={-{Stealth[length=1.6mm]}},
  dim/.style={{Bar[width=2mm]}-{Bar[width=2mm]}, thin},
  cs/.style={draw, fill=ln-accent!55}]
  \node[anchor=east] at (-1,3) {CPU};
  \draw[fill=ln-box] (0,0) rectangle (28,6);   \node at (14,3) {T1};
  \draw[cs]          (28,0) rectangle (33,6);
  \draw[fill=white]  (33,0) rectangle (83,6);  \node at (58,3) {T2};
  \draw[cs]          (83,0) rectangle (88,6);
  \draw[fill=ln-box] (88,0) rectangle (110,6); \node at (99,3) {T1};
  % events
  \draw[ev] (28,13) node[above] {\iflangja{T1 が I/O でブロック}{T1 blocks on I/O}} -- (28,7);
  \draw[ev] (83,13) node[above] {\iflangja{I/O 完了}{I/O completes}} -- (83,7);
  % dimensions
  \draw[dim] (28,-3) -- (33,-3);
  \draw[dim] (83,-3) -- (88,-3);
  \node[below] at (30.5,-3.5) {$t_{\mathrm{ctx}}$};
  \node[below] at (85.5,-3.5) {$t_{\mathrm{ctx}}$};
  \draw[dim] (28,-11) -- (83,-11);
  \node[below] at (55.5,-11.5) {$t_{\mathrm{idle}}$};
\end{tikzpicture}
   (width ~112 mm)
\begin{tikzpicture}[x=1mm, y=1mm, font=\footnotesize\sffamily,
  ev/.style={-{Stealth[length=1.6mm]}},
  dim/.style={{Bar[width=2mm]}-{Bar[width=2mm]}, thin},
  cs/.style={draw, fill=ln-accent!55}]
  \node[anchor=east] at (-1,3) {CPU};
  \draw[fill=ln-box] (0,0) rectangle (28,6);   \node at (14,3) {T1};
  \draw[cs]          (28,0) rectangle (33,6);
  \draw[fill=white]  (33,0) rectangle (83,6);  \node at (58,3) {T2};
  \draw[cs]          (83,0) rectangle (88,6);
  \draw[fill=ln-box] (88,0) rectangle (110,6); \node at (99,3) {T1};
  % events
  \draw[ev] (28,13) node[above] {\iflangja{T1 が I/O でブロック}{T1 blocks on I/O}} -- (28,7);
  \draw[ev] (83,13) node[above] {\iflangja{I/O 完了}{I/O completes}} -- (83,7);
  % dimensions
  \draw[dim] (28,-3) -- (33,-3);
  \draw[dim] (83,-3) -- (88,-3);
  \node[below] at (30.5,-3.5) {$t_{\mathrm{ctx}}$};
  \node[below] at (85.5,-3.5) {$t_{\mathrm{ctx}}$};
  \draw[dim] (28,-11) -- (83,-11);
  \node[below] at (55.5,-11.5) {$t_{\mathrm{idle}}$};
\end{tikzpicture}

  \caption{T1 blocks on I/O; instead of idling, the CPU switches to T2, and
    the completion of the I/O triggers the switch back.  The switches pay
    off when they cost less than the idle time (\cref{eq:l03-idle}).}
  \label{fig:l03-idle-time}
\end{figure}

\begin{example}
  Assume a thread context switch costs \SI{0.05}{\milli\second} and a
  process context switch \SI{0.5}{\milli\second}.  With threads, every
  blocking wait longer than \SI{0.1}{\milli\second} can be used by
  another thread; with processes only waits longer than
  \SI{1}{\milli\second}.  A wait of \SI{4}{\milli\second} gives another
  thread $4 - 0.05 = \SI{3.95}{\milli\second}$ of CPU time and delays the
  blocked thread by \SI{0.05}{\milli\second}; for another process the
  figures are \SI{3.5}{\milli\second} and
  \SI{0.5}{\milli\second}.\margin{The round figures keep the arithmetic
  simple; measured costs are microseconds. lmbench reported
  \SI{6}{\micro\second} per context switch on a 200-MHz P6 in 1996, and
  2--3 GHz cores stay below a microsecond (OSTEP ch.~6).}
\end{example}

\subsection{Multithreaded applications and kernels}

Applications and operating systems both benefit.  A multithreaded
application achieves parallel speedup on several CPUs and hides idle time on one.  A multithreaded
operating system kernel runs some of its threads on behalf of applications
and others to provide operating-system services such as daemons and device
drivers; on a multiprocessor these kernel threads run concurrently on
different CPUs.

\begin{keypoint}
  Threads share the address space of their process, which makes them
  cheaper than processes in memory, context-switch time and communication.
  The same sharing means that threads can interfere with one another, so
  they need explicit mechanisms for coordination.
\end{keypoint}

\section{Basic thread mechanisms}
\label{sec:l03-mechanisms}

Supporting threads requires three things: a data structure that represents
a thread, so that threads can be identified and their resource usage
tracked; mechanisms to create and manage threads; and mechanisms that let
threads running at the same time in the same address space coordinate
safely.\source{P2L2, basic thread mechanisms, thread creation; Birrell
\cite{birrell1989}; OSTEP ch.~26--27.}

Two threads execute concurrently when their executions overlap in time,
whether they are interleaved on one CPU or run simultaneously on different
CPUs; the resulting \term{concurrency}[並行性] creates a problem that
processes do not have.  The operating system gives different processes
different mappings, so a process cannot touch the memory of another.
Threads share their mappings, so any thread can legally read and write any
data of its process.  If two threads update the same data at the same time,
or one reads while another modifies it, the data can become inconsistent.

\begin{definition}[Race condition]
  A \term{race condition}[競合状態] exists when the result of a computation
  depends on the relative timing of threads that access shared data, at
  least one of them modifying it.  Two unsynchronized concurrent accesses to
  the same data, at least one of them a write, form a \emph{data race};
  data races are a common cause of race conditions.
\end{definition}

The mechanisms that threads use to coordinate their access to shared data
and their progress relative to one another are called
\term{synchronization}[同期] mechanisms.  Two basic mechanisms, both
described by Birrell, are used throughout this lesson:
\begin{itemize}
  \item \textbf{mutual exclusion}: only one thread at a time may execute a
    given operation; this is provided by a \emph{mutex}
    (\cref{sec:l03-mutex});
  \item \textbf{waiting}: a thread waits until another thread establishes a
    specific condition, and is woken up by it; this is provided by a
    \emph{condition variable} (\cref{sec:l03-condvar}).
\end{itemize}

\subsection{Thread creation}
\label{sec:l03-creation}

A thread is represented by a thread data structure that contains its
identifier, its program counter, stack pointer and other register values,
its stack, and attributes the thread library uses to manage
it.  In Birrell's notation, \term{thread creation}[スレッド生成] is
performed by the call \texttt{t1 = Fork(proc, args)}, which creates a new
thread, stores its data structure in \texttt{t1}, and
starts the new thread executing \texttt{proc(args)}.  The creating thread
continues with the statement after \texttt{Fork}, so from this point on both
threads run concurrently.  Despite its name, \texttt{Fork} creates a
thread inside the same process; it is unrelated to the Unix \texttt{fork}
of Lesson~2, which creates a new process.

When the child thread finishes, its result must reach the parent.  Either the
child stores it in a well-defined location of the shared address space, or
the parent obtains it with the \term{join}[ジョイン] operation:
\texttt{Join(t1)} blocks the calling thread until thread \texttt{t1} has
completed and returns its result; at that point the resources allocated to
the child are freed.  Apart from this, parent and child are equals: both can
access all resources of the process, and they differ only in their
execution contexts.

\needspace{9\baselineskip}
\begin{example}
  A parent thread and a child thread both insert into a shared list:
  \begin{lstlisting}[language=C, numbers=none]
Thread t1;
Shared_list list;

t1 = Fork(safe_insert, 7);  // child inserts 7
safe_insert(9);             // parent inserts 9
Join(t1);                   // optional here
\end{lstlisting}
  Which insertion happens first depends on how the threads are scheduled,
  so with insertion at the head the list ends up as either $7\to 9$ or
  $9\to 7$.  The \texttt{Join} is not needed for correctness because the
  parent does not use a result of the child.
\end{example}

The example hides a problem.  Inserting element $e$ at the head of a list
takes three steps: create $e$; read the head pointer and store it in
$e$\texttt{.next}; write the address of $e$ into the head pointer.  If two
threads interleave these steps, one insertion can be lost, as the trace in
\cref{tab:l03-race} shows.

\begin{table}[htbp]
  \centering
  \caption{Two unsynchronized insertions into a list that initially holds
    one element, 3.  Both threads read the old head before either writes
    it, so the write in step~6 overwrites that of step~5 and element~7 is
    lost.}
  \label{tab:l03-race}
  \small
  \begin{tabular}{@{}clll@{}}
    \toprule
    Step & Thread T1 (insert 7) & Thread T2 (insert 9) & List \\
    \midrule
    1 & create \texttt{e7}           &                           & $3$ \\
    2 & \texttt{e7.next = head} ($\to 3$) &                      & $3$ \\
    3 &                              & create \texttt{e9}        & $3$ \\
    4 &                              & \texttt{e9.next = head} ($\to 3$) & $3$ \\
    5 & \texttt{head = e7}           &                           & $7\to 3$ \\
    6 &                              & \texttt{head = e9}        & $9\to 3$ \\
    \bottomrule
  \end{tabular}
\end{table}

\section{Mutual exclusion}
\label{sec:l03-mutex}

The race in \cref{tab:l03-race} disappears if only one thread at a time can
execute the three steps of an insertion.\source{P2L2, mutexes; Birrell
\cite{birrell1989}; OSTEP ch.~28.}  With respect to the other threads, the
three steps must then behave like one indivisible operation: no other
thread can observe or modify the list between them.  This requirement is
called \term{mutual exclusion}[相互排除], and the basic mechanism that
enforces it is the mutex.

\begin{definition}[Mutex]
  A \term{mutex}[ミューテックス] (mutual exclusion lock) is a synchronization
  construct with two operations, lock and unlock.  A thread that locks a
  free mutex acquires it and becomes its owner; a thread that tries to lock
  a mutex owned by another thread blocks until the mutex is released.
\end{definition}

A mutex is represented by a data structure that records at least whether it
is locked, which thread owns it, and the list of threads blocked on it.  The
code executed while a mutex is held is the \term{critical section}[クリティカルセクション]: the
part of the program that performs an operation which only one thread at a
time may perform, typically an update of shared state.

\needspace{6\baselineskip}
Birrell's notation makes the critical section a block:
\begin{lstlisting}[language=C, numbers=none]
Lock(m) {
  // critical section
} // unlock
\end{lstlisting}
The mutex is released implicitly when the block ends.  Most thread
interfaces, including the POSIX threads of Lesson~4, instead have separate
calls \texttt{lock(m)} and \texttt{unlock(m)} around the critical section,
and the programmer must make sure that every path through the critical
section ends with an unlock.\caution{A mutex is not recursive: a thread that
locks a mutex it already holds deadlocks against itself.  POSIX prescribes
this for \texttt{PTHREAD\_MUTEX\_NORMAL}, the behaviour of the default type
on Linux; a recursive mutex has to be requested explicitly.}

\needspace{12\baselineskip}
\begin{example}
  With a mutex that protects the list, the insertion of
  \cref{sec:l03-creation} becomes safe:
  \begin{lstlisting}[language=C, numbers=none]
list<int> my_list;
Mutex m;

void safe_insert(int i) {
  Lock(m) {
    my_list.insert(i);
  } // unlock
}
\end{lstlisting}
  If the parent locks \texttt{m} first, the child blocks in \texttt{Lock(m)}
  until the parent's insertion is complete, and then performs its own.  In
  the scenario of \cref{tab:l03-race} the list ends as $7\to 9\to 3$ or
  $9\to 7\to 3$; an element is never lost.
\end{example}

When the owner releases a mutex on which several threads are blocked, only
one of them acquires it; the others stay blocked.  The specification of a
mutex does not say which one: threads are not guaranteed to acquire the
mutex in the order in which they requested it, and a thread that calls
\texttt{Lock} just as the mutex is released may acquire it ahead of threads
that have been waiting for some time.  Program correctness must
therefore not depend on the order in which waiting threads acquire a
mutex.

\section{Condition variables}
\label{sec:l03-condvar}

Mutual exclusion alone is not enough when a thread must wait for a
condition that other threads establish.\source{P2L2, producer/consumer
example, condition variable, condition variable API; Birrell
\cite{birrell1989}; OSTEP ch.~30.}  A thread that holds a mutex cannot
simply stop and wait for the condition, because while it holds the mutex
no other thread can change the shared state.  The standard illustration is
the
\term{producer-consumer problem}[生産者・消費者問題]: several producer
threads insert items into a shared list, and a consumer thread must print
and remove all items once the list is full.

With only a mutex, the consumer has to lock the mutex, check whether the
list is full, and if it is not, release the mutex and try again later.  This
repeated checking wastes CPU time, because most checks find the list not yet
full.  It would be better if the consumer could sleep and be woken up by the
producer whose insertion fills the list.

\begin{definition}[Condition variable]
  A \term{condition variable}[条件変数] is a synchronization construct,
  used together with a mutex, on which threads wait until some condition on
  the shared state holds, and through which other threads notify them when
  it may hold.
\end{definition}

\needspace{8\baselineskip}
With a condition variable \texttt{list\_full}, the two sides become:
\begin{lstlisting}[language=C, numbers=none]
// consumer: print_and_clear
Lock(m) {
  while (my_list.not_full())
    Wait(m, list_full);
  my_list.print_and_remove_all();
} // unlock
\end{lstlisting}
\needspace{7\baselineskip}
\begin{lstlisting}[language=C, numbers=none]
// producers: safe_insert
Lock(m) {
  my_list.insert(my_thread_id);
  if (my_list.full())
    Signal(list_full);
} // unlock
\end{lstlisting}

\subsection{The condition variable interface}

Birrell's interface has three operations on a condition variable
\texttt{cond}.
\begin{description}
  \item[\texttt{Wait(mutex, cond)}] The \term{wait operation}[wait操作] is
    called with \texttt{mutex} held.  It releases the mutex and suspends the
    thread on the wait queue of \texttt{cond}; when the thread is woken up,
    the operation re-acquires the mutex before it returns.
  \item[\texttt{Signal(cond)}] The \term{signal operation}[signal操作] wakes
    up one thread waiting on \texttt{cond}, if there is one.
  \item[\texttt{Broadcast(cond)}] The
    \term{broadcast operation}[broadcast操作] wakes up all threads waiting
    on \texttt{cond}.
\end{description}

\needspace{11\baselineskip}
A condition variable is represented by a data structure that holds a
reference to its mutex and the list of waiting threads.  The wait
operation proceeds as follows:
\begin{lstlisting}[language=C, numbers=none]
Wait(mutex, cond) {
  // atomically release mutex and place the
  //   thread on the wait queue of cond
  // ... wait ...
  // remove the thread from the wait queue
  // re-acquire mutex
  // return
}
\end{lstlisting}
Releasing the mutex while waiting is essential: otherwise no other thread
could enter a critical section to change the state and signal.
Re-acquiring it before returning guarantees that the waiting thread again
holds the mutex when it re-examines the shared state.  As a consequence, a
woken thread cannot continue until it holds the mutex again; if the mutex is
still held, for instance by the signalling thread, the woken thread moves
from the wait queue of the condition variable to that of the mutex
(\cref{fig:l03-cv-queues}).  After a
broadcast all woken threads compete for the mutex, so they leave the wait
operation one at a time.\margin{The signal operation on a condition
variable is unrelated to the Unix signals of Lesson~1.}

\begin{figure}[htbp]
  \centering
  % Life of a thread around a mutex m and a condition variable c.
% Usage: % Life of a thread around a mutex m and a condition variable c.
% Usage: % Life of a thread around a mutex m and a condition variable c.
% Usage: \input{figures/l03-cv-queues}   (width ~116 mm)
\begin{tikzpicture}[x=1mm, y=1mm, font=\footnotesize\sffamily,
  st/.style={draw, rounded corners=2pt, fill=ln-box, align=center,
             minimum height=11mm, inner sep=2pt},
  >={Stealth[length=1.8mm]}, lab/.style={font=\scriptsize\sffamily, align=center}]
  \node[st, minimum width=26mm, text width=24mm] (mq) at (13,0)
    {\iflangja{m の待ちキュー}{wait queue of mutex m}};
  \node[st, minimum width=28mm, text width=26mm, fill=white, thick] (cs) at (58,0)
    {\iflangja{クリティカルセクション\\（m を保持）}{critical section\\(owns m)}};
  \node[st, minimum width=26mm, text width=24mm] (cq) at (103,0)
    {\iflangja{c の待ちキュー}{wait queue of condition c}};
  \node[font=\footnotesize\ttfamily, anchor=south] (lk) at (35,21) {Lock(m)};
  \draw (lk.south) -- (35,16);
  \draw[->] (35,16) -| node[lab, pos=0.4, above] {\iflangja{m は使用中}{m held}} (mq.north);
  \draw[->] (35,16) -| node[lab, pos=0.5, above] {\iflangja{m は空き}{m free}} (cs.north);
  \draw[->] (mq) -- node[lab, above] {\iflangja{m を獲得}{acquire m}} (cs);
  \draw[->] (cs) -- node[lab, above] {\texttt{Wait(m,c)}}
                    node[lab, below] {\iflangja{m を解放}{releases m}} (cq);
  \draw[->] ([xshift=8mm]cs.north) |- node[lab, pos=0.75, above] {\texttt{Unlock(m)}} (94,16);
  % signalled: re-acquire m directly if it is free, otherwise queue on m
  \draw (cq.south) -- node[lab, left, pos=0.55] {\texttt{Signal(c)} / \texttt{Broadcast(c)}} (103,-15);
  \draw[->] (103,-15) -| node[lab, pos=0.75, left] {\iflangja{m は空き}{m free}} ([xshift=8mm]cs.south);
  \draw[->] (103,-15) -- (103,-21) -| node[lab, pos=0.25, below] {\iflangja{m は使用中}{m held}} (mq.south);
\end{tikzpicture}
   (width ~116 mm)
\begin{tikzpicture}[x=1mm, y=1mm, font=\footnotesize\sffamily,
  st/.style={draw, rounded corners=2pt, fill=ln-box, align=center,
             minimum height=11mm, inner sep=2pt},
  >={Stealth[length=1.8mm]}, lab/.style={font=\scriptsize\sffamily, align=center}]
  \node[st, minimum width=26mm, text width=24mm] (mq) at (13,0)
    {\iflangja{m の待ちキュー}{wait queue of mutex m}};
  \node[st, minimum width=28mm, text width=26mm, fill=white, thick] (cs) at (58,0)
    {\iflangja{クリティカルセクション\\（m を保持）}{critical section\\(owns m)}};
  \node[st, minimum width=26mm, text width=24mm] (cq) at (103,0)
    {\iflangja{c の待ちキュー}{wait queue of condition c}};
  \node[font=\footnotesize\ttfamily, anchor=south] (lk) at (35,21) {Lock(m)};
  \draw (lk.south) -- (35,16);
  \draw[->] (35,16) -| node[lab, pos=0.4, above] {\iflangja{m は使用中}{m held}} (mq.north);
  \draw[->] (35,16) -| node[lab, pos=0.5, above] {\iflangja{m は空き}{m free}} (cs.north);
  \draw[->] (mq) -- node[lab, above] {\iflangja{m を獲得}{acquire m}} (cs);
  \draw[->] (cs) -- node[lab, above] {\texttt{Wait(m,c)}}
                    node[lab, below] {\iflangja{m を解放}{releases m}} (cq);
  \draw[->] ([xshift=8mm]cs.north) |- node[lab, pos=0.75, above] {\texttt{Unlock(m)}} (94,16);
  % signalled: re-acquire m directly if it is free, otherwise queue on m
  \draw (cq.south) -- node[lab, left, pos=0.55] {\texttt{Signal(c)} / \texttt{Broadcast(c)}} (103,-15);
  \draw[->] (103,-15) -| node[lab, pos=0.75, left] {\iflangja{m は空き}{m free}} ([xshift=8mm]cs.south);
  \draw[->] (103,-15) -- (103,-21) -| node[lab, pos=0.25, below] {\iflangja{m は使用中}{m held}} (mq.south);
\end{tikzpicture}
   (width ~116 mm)
\begin{tikzpicture}[x=1mm, y=1mm, font=\footnotesize\sffamily,
  st/.style={draw, rounded corners=2pt, fill=ln-box, align=center,
             minimum height=11mm, inner sep=2pt},
  >={Stealth[length=1.8mm]}, lab/.style={font=\scriptsize\sffamily, align=center}]
  \node[st, minimum width=26mm, text width=24mm] (mq) at (13,0)
    {\iflangja{m の待ちキュー}{wait queue of mutex m}};
  \node[st, minimum width=28mm, text width=26mm, fill=white, thick] (cs) at (58,0)
    {\iflangja{クリティカルセクション\\（m を保持）}{critical section\\(owns m)}};
  \node[st, minimum width=26mm, text width=24mm] (cq) at (103,0)
    {\iflangja{c の待ちキュー}{wait queue of condition c}};
  \node[font=\footnotesize\ttfamily, anchor=south] (lk) at (35,21) {Lock(m)};
  \draw (lk.south) -- (35,16);
  \draw[->] (35,16) -| node[lab, pos=0.4, above] {\iflangja{m は使用中}{m held}} (mq.north);
  \draw[->] (35,16) -| node[lab, pos=0.5, above] {\iflangja{m は空き}{m free}} (cs.north);
  \draw[->] (mq) -- node[lab, above] {\iflangja{m を獲得}{acquire m}} (cs);
  \draw[->] (cs) -- node[lab, above] {\texttt{Wait(m,c)}}
                    node[lab, below] {\iflangja{m を解放}{releases m}} (cq);
  \draw[->] ([xshift=8mm]cs.north) |- node[lab, pos=0.75, above] {\texttt{Unlock(m)}} (94,16);
  % signalled: re-acquire m directly if it is free, otherwise queue on m
  \draw (cq.south) -- node[lab, left, pos=0.55] {\texttt{Signal(c)} / \texttt{Broadcast(c)}} (103,-15);
  \draw[->] (103,-15) -| node[lab, pos=0.75, left] {\iflangja{m は空き}{m free}} ([xshift=8mm]cs.south);
  \draw[->] (103,-15) -- (103,-21) -| node[lab, pos=0.25, below] {\iflangja{m は使用中}{m held}} (mq.south);
\end{tikzpicture}

  \caption{A thread waiting on condition \texttt{c} releases mutex
    \texttt{m}.  When it is signalled it must first re-acquire \texttt{m};
    if \texttt{m} is held, it moves to the wait queue of the mutex.}
  \label{fig:l03-cv-queues}
\end{figure}

This is why the consumer tests the list in a \texttt{while} loop and not in
an \texttt{if}.  Between the \texttt{Signal} and the moment the consumer
re-acquires the mutex, other threads may acquire it first and change the
list---for example a second consumer that empties it.  Signalling a
condition variable only says that the condition may hold; the woken thread
cannot be guaranteed access to the mutex immediately, and it must re-check
the condition once it has.

\needspace{7\baselineskip}
\begin{keypoint}
  Always call \texttt{Wait} with the mutex held, inside a \texttt{while}
  loop that re-tests the condition.  Signal or broadcast after changing the
  state that the condition depends on.
\end{keypoint}

\section{The readers-writers problem}
\label{sec:l03-rw}

In the \term{readers-writers problem}[リーダ・ライタ問題] some threads, the
readers, only read a shared resource such as a file, and others, the
writers, modify it.\source{P2L2, readers/writer problem and example,
critical section structure (with proxy); OSTEP ch.~31.}  Any
number of readers may access the resource at the same time, but a writer
needs exclusive access: at most one writer, and no readers while it writes.

Protecting the resource with a single mutex is correct but too restrictive,
because it also excludes readers from each other.  The access rules depend
on how many readers and writers are active:
\begin{itemize}
  \item no active readers and no writer: a reader or a writer may enter;
  \item one or more active readers: another reader may enter, a writer
    may not;
  \item an active writer: neither readers nor writers may enter.
\end{itemize}
These rules can be expressed with a single counter: the resource is
\emph{free} when \texttt{resource\_counter} $=0$, being \emph{read} by $n$
readers when \texttt{resource\_counter} $=n>0$, and being \emph{written}
when \texttt{resource\_counter} $=-1$ (\cref{fig:l03-rw-states}).

\begin{figure}[htbp]
  \centering
  % States of the shared resource in the readers-writers solution.
% Usage: % States of the shared resource in the readers-writers solution.
% Usage: % States of the shared resource in the readers-writers solution.
% Usage: \input{figures/l03-rw-states}   (width ~116 mm)
\begin{tikzpicture}[x=1mm, y=1mm, font=\footnotesize\sffamily,
  st/.style={draw, rounded corners=5pt, fill=ln-box, align=center,
             minimum width=26mm, minimum height=11mm, inner sep=2pt},
  >={Stealth[length=1.8mm]}, lab/.style={font=\scriptsize\sffamily, align=center}]
  \node[st] (wr) at (13,0)  {\iflangja{書き込み中}{writing}\\ \texttt{counter} $=-1$};
  \node[st, fill=white] (fr) at (58,0)  {\iflangja{空き}{free}\\ \texttt{counter} $=0$};
  \node[st] (rd) at (103,0) {\iflangja{読み出し中}{reading}\\ \texttt{counter} $=n>0$};
  \draw[->] (fr.north west) to[bend right=30]
    node[lab, above] {\iflangja{ライタが入る}{writer enters}} (wr.north east);
  \draw[->] (wr.south east) to[bend right=30]
    node[lab, below] {\iflangja{ライタが出る}{writer exits}} (fr.south west);
  \draw[->] (fr.north east) to[bend left=30]
    node[lab, above] {\iflangja{リーダが入る}{reader enters}} (rd.north west);
  \draw[->] (rd.south west) to[bend left=30]
    node[lab, below] {\iflangja{最後のリーダが出る}{last reader exits}} (fr.south east);
  \path[->] (rd) edge[loop above, min distance=7mm, in=60, out=120]
    node[lab, above] {\iflangja{他のリーダが入る/出る}{another reader enters or exits}} (rd);
\end{tikzpicture}
   (width ~116 mm)
\begin{tikzpicture}[x=1mm, y=1mm, font=\footnotesize\sffamily,
  st/.style={draw, rounded corners=5pt, fill=ln-box, align=center,
             minimum width=26mm, minimum height=11mm, inner sep=2pt},
  >={Stealth[length=1.8mm]}, lab/.style={font=\scriptsize\sffamily, align=center}]
  \node[st] (wr) at (13,0)  {\iflangja{書き込み中}{writing}\\ \texttt{counter} $=-1$};
  \node[st, fill=white] (fr) at (58,0)  {\iflangja{空き}{free}\\ \texttt{counter} $=0$};
  \node[st] (rd) at (103,0) {\iflangja{読み出し中}{reading}\\ \texttt{counter} $=n>0$};
  \draw[->] (fr.north west) to[bend right=30]
    node[lab, above] {\iflangja{ライタが入る}{writer enters}} (wr.north east);
  \draw[->] (wr.south east) to[bend right=30]
    node[lab, below] {\iflangja{ライタが出る}{writer exits}} (fr.south west);
  \draw[->] (fr.north east) to[bend left=30]
    node[lab, above] {\iflangja{リーダが入る}{reader enters}} (rd.north west);
  \draw[->] (rd.south west) to[bend left=30]
    node[lab, below] {\iflangja{最後のリーダが出る}{last reader exits}} (fr.south east);
  \path[->] (rd) edge[loop above, min distance=7mm, in=60, out=120]
    node[lab, above] {\iflangja{他のリーダが入る/出る}{another reader enters or exits}} (rd);
\end{tikzpicture}
   (width ~116 mm)
\begin{tikzpicture}[x=1mm, y=1mm, font=\footnotesize\sffamily,
  st/.style={draw, rounded corners=5pt, fill=ln-box, align=center,
             minimum width=26mm, minimum height=11mm, inner sep=2pt},
  >={Stealth[length=1.8mm]}, lab/.style={font=\scriptsize\sffamily, align=center}]
  \node[st] (wr) at (13,0)  {\iflangja{書き込み中}{writing}\\ \texttt{counter} $=-1$};
  \node[st, fill=white] (fr) at (58,0)  {\iflangja{空き}{free}\\ \texttt{counter} $=0$};
  \node[st] (rd) at (103,0) {\iflangja{読み出し中}{reading}\\ \texttt{counter} $=n>0$};
  \draw[->] (fr.north west) to[bend right=30]
    node[lab, above] {\iflangja{ライタが入る}{writer enters}} (wr.north east);
  \draw[->] (wr.south east) to[bend right=30]
    node[lab, below] {\iflangja{ライタが出る}{writer exits}} (fr.south west);
  \draw[->] (fr.north east) to[bend left=30]
    node[lab, above] {\iflangja{リーダが入る}{reader enters}} (rd.north west);
  \draw[->] (rd.south west) to[bend left=30]
    node[lab, below] {\iflangja{最後のリーダが出る}{last reader exits}} (fr.south east);
  \path[->] (rd) edge[loop above, min distance=7mm, in=60, out=120]
    node[lab, above] {\iflangja{他のリーダが入る/出る}{another reader enters or exits}} (rd);
\end{tikzpicture}

  \caption{The states of the shared resource and the corresponding values
    of the proxy variable; \texttt{counter} is short for
    \texttt{resource\_counter}.}
  \label{fig:l03-rw-states}
\end{figure}

The counter is an example of a \term{proxy variable}[代理変数]: a shared
variable, protected by a mutex, that encodes the state of the resource.
Threads hold the mutex only while they examine and update the proxy
variable; the actual reading and writing happens outside the mutex, under
the policy that the proxy variable enforces.  The solution uses one mutex
and two condition variables:
\begin{lstlisting}[language=C, numbers=none]
Mutex counter_mutex;
Condition read_phase, write_phase;
int resource_counter = 0;

// READERS
Lock(counter_mutex) {
  while (resource_counter == -1)
    Wait(counter_mutex, read_phase);
  resource_counter++;
} // unlock
// ... read data ...
Lock(counter_mutex) {
  resource_counter--;
  if (resource_counter == 0)
    Signal(write_phase);
} // unlock
\end{lstlisting}
\needspace{14\baselineskip}
\begin{lstlisting}[language=C, numbers=none]
// WRITER
Lock(counter_mutex) {
  while (resource_counter != 0)
    Wait(counter_mutex, write_phase);
  resource_counter = -1;
} // unlock
// ... write data ...
Lock(counter_mutex) {
  resource_counter = 0;
  Broadcast(read_phase);
  Signal(write_phase);
} // unlock
\end{lstlisting}

The choice of notification in the exit code follows from who can be
waiting.\tip{Signal when at most one waiting thread can proceed and any of
them will do; broadcast when the change can admit several waiters, or when
threads waiting on one condition variable wait for different predicates.}
\begin{itemize}
  \item A \textbf{reader} that leaves while other readers remain changes
    nothing that a waiting thread cares about.  When the last reader leaves,
    only writers can be waiting---readers wait only while a writer is
    active---and only one of them may enter, so a \texttt{Signal} on
    \texttt{write\_phase} suffices.
  \item A \textbf{writer} that leaves frees the resource for either kind of
    thread.  All waiting readers may enter together, hence
    \texttt{Broadcast(read\_phase)}; only one writer may enter, hence
    \texttt{Signal(write\_phase)}.  Which thread actually gets in first
    depends on which one acquires \texttt{counter\_mutex} first.  If a
    reader gets in first, the other woken readers enter as well, and the
    writer finds the counter $\neq 0$ in its \texttt{while} loop and waits
    again; if the writer gets in first, the readers find the counter
    $=-1$ and wait again.\caution{The policy favours readers: while new
    readers keep arriving the counter never returns to 0, and a waiting
    writer never enters.  glibc's \texttt{pthread\_rwlock} is
    reader-preferring by default too, and offers a writer-preferring kind
    (Lesson~10).}
\end{itemize}

\begin{example}
  \Cref{tab:l03-rw-trace} traces two readers, a writer that arrives while
  they are reading, and a third reader that arrives while the writer
  writes.
\end{example}

\begin{table}[htbp]
  \centering
  \caption{A trace of the readers-writers solution, with the value of the
    proxy variable after each step.}
  \label{tab:l03-rw-trace}
  \small
  \begin{tabular}{@{}clr@{}}
    \toprule
    Step & Event & Counter \\
    \midrule
    1 & R1 enters                                               & $1$ \\
    2 & R2 enters                                               & $2$ \\
    3 & W1 finds counter $\neq 0$, waits on \texttt{write\_phase} & $2$ \\
    4 & R1 exits; counter $\neq 0$, no notification             & $1$ \\
    5 & R2 exits; counter $=0$, signals \texttt{write\_phase}   & $0$ \\
    6 & W1 re-acquires the mutex, re-tests, enters             & $-1$ \\
    7 & R3 finds counter $=-1$, waits on \texttt{read\_phase}    & $-1$ \\
    8 & W1 exits; broadcasts \texttt{read\_phase}, signals \texttt{write\_phase} & $0$ \\
    9 & R3 re-acquires the mutex, re-tests, enters             & $1$ \\
    \bottomrule
  \end{tabular}
\end{table}

\subsection{Structure of critical sections}

The readers-writers solution shows the general shape of code that uses a
mutex together with condition variables.  A critical section that may have to wait is written as
\needspace{8\baselineskip}
\begin{lstlisting}[language=C, numbers=none]
Lock(mutex) {
  while (!predicate_indicating_access_ok)
    Wait(mutex, cond_var);
  update state => update predicate;
  Signal and/or Broadcast(
      cond_var_with_correct_waiting_threads);
} // unlock
\end{lstlisting}
\needspace{9\baselineskip}
With a proxy variable, the operation on the resource is split out of the
mutex into three parts:
\begin{lstlisting}[language=C, numbers=none]
// ENTER CRITICAL SECTION
Lock(mutex) {
  while (!predicate_for_access)
    Wait(mutex, cond_var);
  update predicate;
} // unlock

perform critical operation  // e.g. read/write shared file

// EXIT CRITICAL SECTION
Lock(mutex) {
  update predicate;
  Signal/Broadcast(cond_var);
} // unlock
\end{lstlisting}

\needspace{6\baselineskip}
\begin{keypoint}
  A mutex by itself enforces one policy: one thread at a time.  Proxy
  variables, updated in short critical sections, allow more complex access
  policies---such as \enquote{any number of readers or one writer}---to be
  built from the same mutex and condition variable primitives.
\end{keypoint}

\section{Common pitfalls and spurious wake-ups}
\label{sec:l03-pitfalls}

Several mistakes recur in programs that use mutexes and condition
variables.\source{P2L2, common pitfalls, spurious wake-ups; Birrell
\cite{birrell1989}.}
\begin{itemize}
  \item Losing track of which mutex and condition variables belong to which
    resource.  Recording the association where the variables are declared,
    e.g.\ \texttt{Mutex m1; // protects file1}, avoids this.
  \item Forgetting to lock or to unlock, or locking and unlocking the wrong
    mutex.  Some compilers and analysis tools can warn about such
    cases.\margin{Clang's \texttt{-Wthread-safety} checks lock annotations
    while compiling; ThreadSanitizer (\texttt{-fsanitize=thread}) finds races
    while the program runs, at a typical cost of 5--15$\times$ in time and
    5--10$\times$ in memory.}
  \item Using more than one mutex for the same resource.  If the readers of
    a file lock \texttt{m1} and its writers lock \texttt{m2}, reads and
    writes run concurrently.
  \item Signalling the wrong condition variable.
  \item Using \texttt{Signal} where \texttt{Broadcast} is needed.  Only one
    thread is woken; the others may continue to wait indefinitely.
  \item Assuming that notifications establish an execution order.  The
    order in which signals are issued does not determine the order in
    which the woken threads run; if priorities must be respected, another
    mechanism is needed.
\end{itemize}

Two further problems deserve a closer look: spurious wake-ups, which cost
performance, and deadlocks (\cref{sec:l03-deadlock}), which stop a program.

\subsection{Spurious wake-ups}

Some notifications wake threads that cannot use them.

\begin{definition}[Spurious wake-up]
  A \term{spurious wake-up}[スプリアスウェイクアップ] is a wake-up of a
  thread that cannot make progress when it runs, so that the context switches
  it causes are wasted.
\end{definition}

A common cause is signalling while the mutex is still held.  In the writer's exit code of \cref{sec:l03-rw} the
broadcast is issued inside the critical section.  The woken readers leave the
wait queue of \texttt{read\_phase}, but they cannot return from
\texttt{Wait} before re-acquiring \texttt{counter\_mutex}, which the writer
still holds; they are simply moved to the wait queue of the mutex
(\cref{fig:l03-cv-queues}).\margin{POSIX also calls a return from wait
without any signal a spurious wake-up; the \texttt{while} loop handles
both.}

\needspace{9\baselineskip}
When the notifications do not depend on the shared state, the fix is to
move them after the unlock:
\begin{lstlisting}[language=C, numbers=none]
// writer exit, notifications outside the mutex
Lock(counter_mutex) {
  resource_counter = 0;
} // unlock
Broadcast(read_phase);
Signal(write_phase);
\end{lstlisting}

This is not always possible.  In the reader's exit code the decision whether
to signal depends on \texttt{resource\_counter}, which may only be read while
holding the mutex, so the \texttt{Signal} cannot simply be moved after the
unlock; it can be moved only by recording the decision in a local variable
inside the critical section and signalling afterwards.\margin{Linux offers
\texttt{FUTEX\_CMP\_REQUEUE} for the same purpose: instead of waking the
threads waiting on a condition variable it moves them to the wait queue of
the mutex, which \texttt{futex(2)} describes as avoiding a
\enquote{thundering herd}.}  Spurious wake-ups
never affect correctness---the \texttt{while} loop sends a thread that cannot
proceed back to waiting---only performance.

\section{Deadlocks}
\label{sec:l03-deadlock}

Threads that hold some mutexes while requesting others can block each other
permanently.\source{P2L2, deadlocks; OSTEP ch.~32; Silberschatz et al.\
ch.~8.}

\begin{definition}[Deadlock]
  A \term{deadlock}[デッドロック] is a situation in which two or more
  competing threads each wait for another to complete, but none of them
  ever does.
\end{definition}

\needspace{10\baselineskip}
The simplest case involves two threads and two mutexes.  Both threads need resources A and B to execute
\texttt{foo(A, B)}, but they lock the protecting mutexes in opposite
order:
\begin{lstlisting}[language=C, numbers=none]
// thread T1              // thread T2
Lock(m_A) {               Lock(m_B) {
  Lock(m_B) {               Lock(m_A) {
    foo(A, B);                foo(A, B);
  }                         }
}                         }
\end{lstlisting}
If T1 acquires \texttt{m\_A} and T2 acquires \texttt{m\_B} before either
reaches its second \texttt{Lock}, T1 blocks waiting for \texttt{m\_B} and T2
blocks waiting for \texttt{m\_A}.  Neither can proceed, and neither will ever
release the mutex the other needs.

The situation is described by a \term{wait graph}[待ちグラフ], which has an
edge from a thread waiting for a resource to the thread that owns the
resource.  \Cref{fig:l03-wait-graph} shows the example with the mutexes drawn
explicitly.  A deadlock corresponds to a cycle in the wait graph: for locks,
which have a single owner, a cycle is both necessary and sufficient for a
deadlock.

\begin{figure}[htbp]
  \centering
  % Wait graph of the two-thread, two-lock deadlock.
% Usage: % Wait graph of the two-thread, two-lock deadlock.
% Usage: % Wait graph of the two-thread, two-lock deadlock.
% Usage: \input{figures/l03-wait-graph}   (width ~80 mm)
\begin{tikzpicture}[x=1mm, y=1mm, font=\footnotesize\sffamily,
  th/.style={draw, circle, fill=ln-box, minimum size=9mm, inner sep=0pt},
  lk/.style={draw, rectangle, fill=white, thick, minimum size=7mm, inner sep=0pt},
  >={Stealth[length=1.8mm]}, lab/.style={font=\scriptsize\sffamily},
  holds/.style={->, thick}, waits/.style={->, thick, dashed, ln-caution}]
  \node[th] (t1) at (0,0)   {T1};
  \node[th] (t2) at (60,0)  {T2};
  \node[lk] (a)  at (30,14) {\texttt{m\_A}};
  \node[lk] (b)  at (30,-14) {\texttt{m\_B}};
  \draw[holds] (a) -- node[lab, above left]  {\iflangja{T1 が保持}{held by T1}} (t1);
  \draw[waits] (t1) -- node[lab, below left] {\iflangja{T1 が待つ}{T1 waits for}} (b);
  \draw[holds] (b) -- node[lab, below right] {\iflangja{T2 が保持}{held by T2}} (t2);
  \draw[waits] (t2) -- node[lab, above right] {\iflangja{T2 が待つ}{T2 waits for}} (a);
\end{tikzpicture}
   (width ~80 mm)
\begin{tikzpicture}[x=1mm, y=1mm, font=\footnotesize\sffamily,
  th/.style={draw, circle, fill=ln-box, minimum size=9mm, inner sep=0pt},
  lk/.style={draw, rectangle, fill=white, thick, minimum size=7mm, inner sep=0pt},
  >={Stealth[length=1.8mm]}, lab/.style={font=\scriptsize\sffamily},
  holds/.style={->, thick}, waits/.style={->, thick, dashed, ln-caution}]
  \node[th] (t1) at (0,0)   {T1};
  \node[th] (t2) at (60,0)  {T2};
  \node[lk] (a)  at (30,14) {\texttt{m\_A}};
  \node[lk] (b)  at (30,-14) {\texttt{m\_B}};
  \draw[holds] (a) -- node[lab, above left]  {\iflangja{T1 が保持}{held by T1}} (t1);
  \draw[waits] (t1) -- node[lab, below left] {\iflangja{T1 が待つ}{T1 waits for}} (b);
  \draw[holds] (b) -- node[lab, below right] {\iflangja{T2 が保持}{held by T2}} (t2);
  \draw[waits] (t2) -- node[lab, above right] {\iflangja{T2 が待つ}{T2 waits for}} (a);
\end{tikzpicture}
   (width ~80 mm)
\begin{tikzpicture}[x=1mm, y=1mm, font=\footnotesize\sffamily,
  th/.style={draw, circle, fill=ln-box, minimum size=9mm, inner sep=0pt},
  lk/.style={draw, rectangle, fill=white, thick, minimum size=7mm, inner sep=0pt},
  >={Stealth[length=1.8mm]}, lab/.style={font=\scriptsize\sffamily},
  holds/.style={->, thick}, waits/.style={->, thick, dashed, ln-caution}]
  \node[th] (t1) at (0,0)   {T1};
  \node[th] (t2) at (60,0)  {T2};
  \node[lk] (a)  at (30,14) {\texttt{m\_A}};
  \node[lk] (b)  at (30,-14) {\texttt{m\_B}};
  \draw[holds] (a) -- node[lab, above left]  {\iflangja{T1 が保持}{held by T1}} (t1);
  \draw[waits] (t1) -- node[lab, below left] {\iflangja{T1 が待つ}{T1 waits for}} (b);
  \draw[holds] (b) -- node[lab, below right] {\iflangja{T2 が保持}{held by T2}} (t2);
  \draw[waits] (t2) -- node[lab, above right] {\iflangja{T2 が待つ}{T2 waits for}} (a);
\end{tikzpicture}

  \caption{Wait graph of the deadlock: following the edges from T1 leads
    through \texttt{m\_B}, T2 and \texttt{m\_A} back to T1.}
  \label{fig:l03-wait-graph}
\end{figure}

\subsection{Avoiding the example deadlock}

There are three ways to remove the deadlock from the example.
\begin{enumerate}
  \item \textbf{Unlock A before locking B.}  Holding only one mutex at a
    time (fine-grained locking) makes a cycle impossible, but it does not
    work here: the threads need both resources at once.
  \item \textbf{Acquire all locks up front.}  A thread that obtains every
    lock it needs at once before it starts---atomically, all or none,
    releasing any it holds if one is unavailable---and releases them all at
    the end cannot deadlock; the extreme case is a single \enquote{mega
    lock} for all resources.  For many applications this is too restrictive, because it
    limits parallelism.
  \item \textbf{Maintain a lock order.}  If every thread acquires the
    mutexes in the same order, first \texttt{m\_A} and then \texttt{m\_B},
    the deadlock cannot occur.
\end{enumerate}

The third approach, \term{lock ordering}[ロック順序], works in general.
Number the mutexes and require every thread to acquire them in increasing
order.  Suppose threads $T_1,\dots,T_k$ formed a cycle in which each $T_j$
waits for a mutex $M_j$ owned by the next thread, $T_{j+1}$ (with
$T_{k+1}=T_1$).  Thread $T_{j+1}$ holds $M_j$ and requests $M_{j+1}$, so
$M_{j+1}>M_j$.  The numbers would then increase strictly all the way
around the cycle and yet return to $M_1$, which is impossible.  The
difficulty is practical: in a large program with many mutexes it is hard to
ensure that every code path respects the order, and tools that check lock
acquisition orders help.

\subsection{Handling deadlocks in general}

Operating systems and applications can deal with deadlocks in three
ways.
\begin{description}
  \item[Prevention] With \term{deadlock prevention}[デッドロック防止] a cycle
    is never allowed to form.  Either the code is changed so that cycles are
    impossible---a thread releases the locks it holds before it requests
    another, or all threads follow a lock order---or every lock request is
    checked at run time and delayed if granting it could later lead to a
    cycle.  The run-time check is expensive, and it requires knowing in
    advance which locks each thread may request; textbooks call this variant
    \emph{deadlock avoidance} and reserve \emph{prevention} for the
    structural rules.
  \item[Detection and recovery] With
    \term{deadlock detection and recovery}[デッドロックの検出と回復] the
    system periodically analyses the wait graph for cycles and, when it
    finds one, rolls back the execution of some of the threads involved.
    Rollback requires enough saved state to restore an earlier point, and
    it is impossible when the threads have already consumed external input
    or produced external output.
  \item[Doing nothing] The \term{ostrich algorithm}[ダチョウアルゴリズム]
    ignores deadlocks altogether, on the grounds that they are rare and the
    other approaches are expensive.  If a deadlock does occur, the system is
    rebooted.
\end{description}

\begin{keypoint}
  A deadlock is a cycle in the wait graph.  Acquiring all mutexes in one
  global order prevents such cycles at no run-time cost; checking every
  lock request at run time and recovery by rollback are general but
  expensive.
\end{keypoint}

\section{Kernel-level and user-level threads}
\label{sec:l03-klt-ult}

Threads exist at two levels of the system.\source{P2L2, kernel vs.\
user-level threads, multithreading models, scope of multithreading;
Silberschatz et al.\ ch.~4--5.}  The operating system kernel can itself be
multithreaded, and applications can create threads within their processes;
the question is how the two kinds of threads relate to each other.
\begin{definition}[Kernel-level and user-level threads]
  A \term{kernel-level thread}[カーネルレベルスレッド] is a thread that is
  visible to the kernel and managed by kernel components such as the CPU
  scheduler.  A \term{user-level thread}[ユーザレベルスレッド] is a thread
  as the application sees it, created through a user-level thread library;
  depending on the multithreading model, the library manages it itself or
  passes the operations to the kernel.
\end{definition}

Some kernel-level threads exist to support the user-level threads of
applications; others run operating-system services such as daemons.  The
CPU scheduler decides which kernel-level threads run on which CPUs.  A
user-level thread, in turn, can execute only when it is associated with a
kernel-level thread and the scheduler runs that kernel-level thread on a CPU.

\subsection{Multithreading models}

How user-level threads are associated with kernel-level threads is
described by three multithreading models
(\cref{fig:l03-mt-models}).\sidenote{The many-to-many model was the
ambitious one, and it has largely been abandoned.  Solaris~2 implemented it,
but the coordination between the two levels proved to cost more than it
saved, and Sun made the one-to-one library the default in Solaris~9 (2002);
Linux settled on the one-to-one NPTL rather than the many-to-many NGPT
(Lesson~5).  The model survives in language runtimes, where the scheduler is
part of the runtime: a Go program multiplexes goroutines, which start with a
stack of \SI{2}{\kibi\byte}, onto a few kernel-level threads.}

\begin{figure}[htbp]
  \centering
  % Multithreading models: mapping of user-level threads to kernel-level threads.
% Usage: % Multithreading models: mapping of user-level threads to kernel-level threads.
% Usage: % Multithreading models: mapping of user-level threads to kernel-level threads.
% Usage: \input{figures/l03-mt-models}   (width ~118 mm)
\begin{tikzpicture}[x=0.96mm, y=1mm, font=\footnotesize\sffamily,
  ut/.style={draw, circle, fill=white, minimum size=4mm, inner sep=0pt},
  kt/.style={draw, rectangle, fill=ln-box, minimum size=4mm, inner sep=0pt},
  lib/.style={draw, rounded corners=1pt, fill=ln-box, minimum height=4mm,
              inner sep=1pt, font=\scriptsize\sffamily},
  ttl/.style={font=\footnotesize\sffamily\bfseries},
  lab/.style={font=\scriptsize\sffamily, text=ln-muted}]
  \draw[dashed, ln-rule] (0,10) -- (118,10);
  \node[lab, anchor=south west] at (0,10.5) {\iflangja{ユーザ}{user}};
  \node[lab, anchor=north west] at (0,9.5)  {\iflangja{カーネル}{kernel}};
  % ---- one-to-one
  \node[ttl] at (27,32) {\iflangja{一対一}{one-to-one}};
  \foreach \x in {17,27,37} {
    \node[ut] (u) at (\x,24) {};
    \node[kt] (k) at (\x,3) {};
    \draw (u) -- (k);
  }
  % ---- many-to-one
  \node[ttl] at (65,32) {\iflangja{多対一}{many-to-one}};
  \node[lib, minimum width=28mm] (l2) at (65,16) {\iflangja{スレッドライブラリ}{thread library}};
  \foreach \x in {55,65,75} { \node[ut] (u) at (\x,24) {}; \draw (u) -- (\x,18); }
  \node[kt] (k2) at (65,3) {};
  \draw (l2) -- (k2);
  % ---- many-to-many
  \node[ttl] at (103,32) {\iflangja{多対多}{many-to-many}};
  \node[lib, minimum width=24mm] (l3) at (99,16) {\iflangja{ライブラリ}{library}};
  \foreach \x in {90,99,108} { \node[ut] (u) at (\x,24) {}; \draw (u) -- (\x,18); }
  \node[kt] (k3a) at (93,3) {};
  \node[kt] (k3b) at (105,3) {};
  \draw (l3.south) ++(-6,0) -- (k3a);
  \draw (l3.south) ++(6,0) -- (k3b);
  \node[ut, thick, draw=ln-accent] (ub) at (116,24) {};
  \node[kt, thick, draw=ln-accent] (kb) at (116,3) {};
  \draw[very thick, ln-accent] (ub) -- (kb);
  \node[lab, anchor=north, text=ln-accent] at (116,0.5) {\iflangja{バウンド}{bound}};
\end{tikzpicture}
   (width ~118 mm)
\begin{tikzpicture}[x=0.96mm, y=1mm, font=\footnotesize\sffamily,
  ut/.style={draw, circle, fill=white, minimum size=4mm, inner sep=0pt},
  kt/.style={draw, rectangle, fill=ln-box, minimum size=4mm, inner sep=0pt},
  lib/.style={draw, rounded corners=1pt, fill=ln-box, minimum height=4mm,
              inner sep=1pt, font=\scriptsize\sffamily},
  ttl/.style={font=\footnotesize\sffamily\bfseries},
  lab/.style={font=\scriptsize\sffamily, text=ln-muted}]
  \draw[dashed, ln-rule] (0,10) -- (118,10);
  \node[lab, anchor=south west] at (0,10.5) {\iflangja{ユーザ}{user}};
  \node[lab, anchor=north west] at (0,9.5)  {\iflangja{カーネル}{kernel}};
  % ---- one-to-one
  \node[ttl] at (27,32) {\iflangja{一対一}{one-to-one}};
  \foreach \x in {17,27,37} {
    \node[ut] (u) at (\x,24) {};
    \node[kt] (k) at (\x,3) {};
    \draw (u) -- (k);
  }
  % ---- many-to-one
  \node[ttl] at (65,32) {\iflangja{多対一}{many-to-one}};
  \node[lib, minimum width=28mm] (l2) at (65,16) {\iflangja{スレッドライブラリ}{thread library}};
  \foreach \x in {55,65,75} { \node[ut] (u) at (\x,24) {}; \draw (u) -- (\x,18); }
  \node[kt] (k2) at (65,3) {};
  \draw (l2) -- (k2);
  % ---- many-to-many
  \node[ttl] at (103,32) {\iflangja{多対多}{many-to-many}};
  \node[lib, minimum width=24mm] (l3) at (99,16) {\iflangja{ライブラリ}{library}};
  \foreach \x in {90,99,108} { \node[ut] (u) at (\x,24) {}; \draw (u) -- (\x,18); }
  \node[kt] (k3a) at (93,3) {};
  \node[kt] (k3b) at (105,3) {};
  \draw (l3.south) ++(-6,0) -- (k3a);
  \draw (l3.south) ++(6,0) -- (k3b);
  \node[ut, thick, draw=ln-accent] (ub) at (116,24) {};
  \node[kt, thick, draw=ln-accent] (kb) at (116,3) {};
  \draw[very thick, ln-accent] (ub) -- (kb);
  \node[lab, anchor=north, text=ln-accent] at (116,0.5) {\iflangja{バウンド}{bound}};
\end{tikzpicture}
   (width ~118 mm)
\begin{tikzpicture}[x=0.96mm, y=1mm, font=\footnotesize\sffamily,
  ut/.style={draw, circle, fill=white, minimum size=4mm, inner sep=0pt},
  kt/.style={draw, rectangle, fill=ln-box, minimum size=4mm, inner sep=0pt},
  lib/.style={draw, rounded corners=1pt, fill=ln-box, minimum height=4mm,
              inner sep=1pt, font=\scriptsize\sffamily},
  ttl/.style={font=\footnotesize\sffamily\bfseries},
  lab/.style={font=\scriptsize\sffamily, text=ln-muted}]
  \draw[dashed, ln-rule] (0,10) -- (118,10);
  \node[lab, anchor=south west] at (0,10.5) {\iflangja{ユーザ}{user}};
  \node[lab, anchor=north west] at (0,9.5)  {\iflangja{カーネル}{kernel}};
  % ---- one-to-one
  \node[ttl] at (27,32) {\iflangja{一対一}{one-to-one}};
  \foreach \x in {17,27,37} {
    \node[ut] (u) at (\x,24) {};
    \node[kt] (k) at (\x,3) {};
    \draw (u) -- (k);
  }
  % ---- many-to-one
  \node[ttl] at (65,32) {\iflangja{多対一}{many-to-one}};
  \node[lib, minimum width=28mm] (l2) at (65,16) {\iflangja{スレッドライブラリ}{thread library}};
  \foreach \x in {55,65,75} { \node[ut] (u) at (\x,24) {}; \draw (u) -- (\x,18); }
  \node[kt] (k2) at (65,3) {};
  \draw (l2) -- (k2);
  % ---- many-to-many
  \node[ttl] at (103,32) {\iflangja{多対多}{many-to-many}};
  \node[lib, minimum width=24mm] (l3) at (99,16) {\iflangja{ライブラリ}{library}};
  \foreach \x in {90,99,108} { \node[ut] (u) at (\x,24) {}; \draw (u) -- (\x,18); }
  \node[kt] (k3a) at (93,3) {};
  \node[kt] (k3b) at (105,3) {};
  \draw (l3.south) ++(-6,0) -- (k3a);
  \draw (l3.south) ++(6,0) -- (k3b);
  \node[ut, thick, draw=ln-accent] (ub) at (116,24) {};
  \node[kt, thick, draw=ln-accent] (kb) at (116,3) {};
  \draw[very thick, ln-accent] (ub) -- (kb);
  \node[lab, anchor=north, text=ln-accent] at (116,0.5) {\iflangja{バウンド}{bound}};
\end{tikzpicture}

  \caption{Multithreading models.  Circles are user-level threads, squares
    kernel-level threads.  In the many-to-many model a bound thread keeps a
    kernel-level thread of its own.}
  \label{fig:l03-mt-models}
\end{figure}

\begin{description}
  \item[One-to-one] In the \term{one-to-one model}[一対一モデル] every
    user-level thread has its own kernel-level thread.  The kernel sees and
    understands the threads, their synchronization and their blocking, and
    can schedule them on several CPUs.  On the other hand every thread
    operation goes through the kernel with a system call, which is
    expensive; the application is limited by the policies and the maximum
    number of threads that the kernel supports; and it is portable only to
    kernels that provide this support.
  \item[Many-to-one] In the \term{many-to-one model}[多対一モデル] all
    user-level threads of a process are mapped onto a single kernel-level
    thread, and a user-level thread library decides which user-level thread
    runs on it at any time.  Thread management needs no system calls and
    does not depend on kernel limits or policies, so it is fully portable.
    But the kernel has no insight into the application's needs, and when
    one user-level thread blocks on I/O, the kernel-level thread blocks, and
    with it the whole process.
  \item[Many-to-many] In the \term{many-to-many model}[多対多モデル] some
    user-level threads are multiplexed onto a set of kernel-level threads
    while others have a kernel-level thread to themselves.  A user-level
    thread that is permanently associated with a kernel-level thread is a
    \term{bound thread}[バウンドスレッド]; it can, for example, be given a
    higher priority or better responsiveness.  Because the process has
    several kernel-level threads, a user-level thread that blocks on I/O
    blocks only its kernel-level thread; the other user-level threads
    continue on the remaining ones.  The model combines the
    advantages of the other two, at the price of coordination between the
    user-level and the kernel-level thread managers.
\end{description}


\subsection{Scope of multithreading}

Thread management happens at two levels, and the level determines the unit
among which resources are shared.
With \term{process scope}[プロセススコープ] the user-level thread library
manages the threads within a single process and competes for resources only
with the other threads of that process; the kernel does not see them
individually.  With \term{system scope}[システムスコープ] threads are
managed system-wide by the operating system's thread managers, in
particular the CPU scheduler, and compete with all threads in the system.

\begin{example}
  Process P runs 6 user-level threads and process Q runs 2.  With process
  scope the kernel sees only the two processes and, if it gives each an
  equal share of the CPU, each thread of P receives
  $\frac{1}{2}\cdot\frac{1}{6} = \frac{1}{12}$ of the CPU while each thread
  of Q receives $\frac{1}{4}$.  With system scope the kernel sees all 8
  threads; dividing the CPU equally among them gives every thread
  $\frac{1}{8}$, so that P as a whole receives $\frac{3}{4}$.
\end{example}

\section{Multithreading patterns}
\label{sec:l03-patterns}

Multithreaded applications are commonly organized according to one of three
patterns: boss-workers, pipeline and layered.\source{P2L2, multithreading
patterns; thread pools: Silberschatz et al.\ ch.~4.}

\subsection{Boss-workers}

In the \term{boss-workers pattern}[ボス・ワーカーパターン] one boss thread
accepts requests and assigns them to worker threads, and each worker
performs an entire task.  Every request passes through the boss, so the
throughput of the system is limited by the boss: if the boss spends time
$t_{\mathrm{boss}}$ on each request, the system can accept at most
$1/t_{\mathrm{boss}}$ requests per unit time, however many workers there
are.  The boss must therefore do as little per request as possible.

\begin{figure}[htbp]
  \centering
  % Boss-workers pattern with a shared work queue.
% Usage: % Boss-workers pattern with a shared work queue.
% Usage: % Boss-workers pattern with a shared work queue.
% Usage: \input{figures/l03-boss-workers}   (width ~106 mm)
\begin{tikzpicture}[x=1mm, y=1mm, font=\footnotesize\sffamily,
  box/.style={draw, rounded corners=2pt, fill=ln-box, align=center, inner sep=1pt},
  cell/.style={draw, fill=white, minimum width=3mm, minimum height=5mm, inner sep=0pt},
  >={Stealth[length=1.6mm]},
  lab/.style={font=\scriptsize\sffamily, text=ln-muted}]
  \node[lab, anchor=south west] at (0,0.5) {\iflangja{要求}{requests}};
  \node[box, minimum width=14mm, minimum height=7mm] (boss) at (26,0) {\iflangja{ボス}{boss}};
  \draw[->] (0,0) -- (boss);
  \foreach \i in {0,...,3} { \node[cell] at (46+3*\i,0) {}; }
  \node[lab, anchor=north] at (50.5,-3) {\iflangja{キュー}{queue}};
  \draw[->] (boss) -- (44.5,0);
  \foreach \y/\n in {8/1,0/2,-8/3} {
    \node[box, minimum width=18mm, minimum height=6mm] (w\n) at (82,\y) {\iflangja{ワーカー}{worker} \n};
    \draw[->] (56.5,0) -- (w\n.west);
    \draw[->] (w\n.east) -- ++(12,0);
  }
  \node[lab, anchor=south west] at (93,8.5) {\iflangja{応答}{responses}};
\end{tikzpicture}
   (width ~106 mm)
\begin{tikzpicture}[x=1mm, y=1mm, font=\footnotesize\sffamily,
  box/.style={draw, rounded corners=2pt, fill=ln-box, align=center, inner sep=1pt},
  cell/.style={draw, fill=white, minimum width=3mm, minimum height=5mm, inner sep=0pt},
  >={Stealth[length=1.6mm]},
  lab/.style={font=\scriptsize\sffamily, text=ln-muted}]
  \node[lab, anchor=south west] at (0,0.5) {\iflangja{要求}{requests}};
  \node[box, minimum width=14mm, minimum height=7mm] (boss) at (26,0) {\iflangja{ボス}{boss}};
  \draw[->] (0,0) -- (boss);
  \foreach \i in {0,...,3} { \node[cell] at (46+3*\i,0) {}; }
  \node[lab, anchor=north] at (50.5,-3) {\iflangja{キュー}{queue}};
  \draw[->] (boss) -- (44.5,0);
  \foreach \y/\n in {8/1,0/2,-8/3} {
    \node[box, minimum width=18mm, minimum height=6mm] (w\n) at (82,\y) {\iflangja{ワーカー}{worker} \n};
    \draw[->] (56.5,0) -- (w\n.west);
    \draw[->] (w\n.east) -- ++(12,0);
  }
  \node[lab, anchor=south west] at (93,8.5) {\iflangja{応答}{responses}};
\end{tikzpicture}
   (width ~106 mm)
\begin{tikzpicture}[x=1mm, y=1mm, font=\footnotesize\sffamily,
  box/.style={draw, rounded corners=2pt, fill=ln-box, align=center, inner sep=1pt},
  cell/.style={draw, fill=white, minimum width=3mm, minimum height=5mm, inner sep=0pt},
  >={Stealth[length=1.6mm]},
  lab/.style={font=\scriptsize\sffamily, text=ln-muted}]
  \node[lab, anchor=south west] at (0,0.5) {\iflangja{要求}{requests}};
  \node[box, minimum width=14mm, minimum height=7mm] (boss) at (26,0) {\iflangja{ボス}{boss}};
  \draw[->] (0,0) -- (boss);
  \foreach \i in {0,...,3} { \node[cell] at (46+3*\i,0) {}; }
  \node[lab, anchor=north] at (50.5,-3) {\iflangja{キュー}{queue}};
  \draw[->] (boss) -- (44.5,0);
  \foreach \y/\n in {8/1,0/2,-8/3} {
    \node[box, minimum width=18mm, minimum height=6mm] (w\n) at (82,\y) {\iflangja{ワーカー}{worker} \n};
    \draw[->] (56.5,0) -- (w\n.west);
    \draw[->] (w\n.east) -- ++(12,0);
  }
  \node[lab, anchor=south west] at (93,8.5) {\iflangja{応答}{responses}};
\end{tikzpicture}

  \caption{Boss-workers with a work queue: the boss only enqueues requests;
    every worker takes a request from the queue and performs the entire
    task.}
  \label{fig:l03-boss-workers}
\end{figure}

The boss can pass work to workers in two ways
(\cref{fig:l03-boss-workers} shows the second).
\begin{itemize}
  \item \textbf{Directly signalling a specific worker.}  The workers need
    no synchronization among themselves, but the boss must keep track of
    what each worker is doing and which one is free.  This increases
    $t_{\mathrm{boss}}$ and lowers throughput.
  \item \textbf{Placing work in a queue,} from which the workers take it,
    in a producer-consumer arrangement.  The boss need not know anything about
    the workers, but the workers (and the boss) must synchronize their
    access to the shared queue.  Because it keeps the boss's work small,
    this option usually gives the better throughput.
\end{itemize}

Workers can be created on demand, when a request arrives, but creating a
thread for every request adds latency.  More commonly the application uses
a \term{thread pool}[スレッドプール], a set of workers created in
advance.
The pool can be static, or dynamic, growing---typically by several threads
at a time rather than one---when the queue gets long and shrinking when
workers are idle.\tip{Size a pool by the work, not by the request rate: a
worker that spends a fraction $b$ of each request blocked keeps a CPU busy
only $1-b$ of the time, so $c$ CPUs need about $c/(1-b)$ workers.}

The pattern is simple, but it requires thread pool management, and it
ignores locality: the boss does not know which worker has just performed a
similar task and would perform the next one with a hot cache.  A variant
therefore specializes the workers for particular types of requests instead
of making all workers equal.  Specialized workers have better locality, and
the application can provide different quality of service by assigning more
threads to more important request types; the cost is load balancing---it is
harder to decide how many threads each type should receive.

\subsection{Pipeline}

In the \term{pipeline pattern}[パイプラインパターン] the task is divided
into a sequence of subtasks, each performed by a thread (or a group of
threads) dedicated to it.  Stages communicate through shared buffers, so
several requests are processed concurrently, each in a different stage
(\cref{fig:l03-pipeline}).  The
throughput is determined by the weakest link, the slowest stage.  The
remedy is to make each stage a thread pool and give slow stages more
threads.  If stage $i$ takes time $t_i$ per request and has $p_i$ threads,
the throughput is
\begin{equation}
  X \;=\; \min_i \frac{p_i}{t_i} .
  \label{eq:l03-pipeline-throughput}
\end{equation}

\begin{figure}[htbp]
  \centering
  % Pipeline pattern: stages connected by buffers; stage 2 is a thread pool.
% Usage: % Pipeline pattern: stages connected by buffers; stage 2 is a thread pool.
% Usage: % Pipeline pattern: stages connected by buffers; stage 2 is a thread pool.
% Usage: \input{figures/l03-pipeline}   (width ~116 mm)
\begin{tikzpicture}[x=1mm, y=1mm, font=\footnotesize\sffamily,
  box/.style={draw, rounded corners=2pt, fill=ln-box, align=center, inner sep=1pt},
  th/.style={draw, circle, fill=white, minimum size=4mm, inner sep=0pt},
  cell/.style={draw, fill=white, minimum width=3mm, minimum height=5mm, inner sep=0pt},
  >={Stealth[length=1.6mm]},
  lab/.style={font=\scriptsize\sffamily, text=ln-muted}]
  \node[box, minimum width=14mm, minimum height=17mm] (s1) at (20,0) {};
  \node[box, minimum width=14mm, minimum height=17mm] (s2) at (60,0) {};
  \node[box, minimum width=14mm, minimum height=17mm] (s3) at (100,0) {};
  \foreach \s/\n in {s1/1,s2/2,s3/3} {
    \node[font=\scriptsize\sffamily, anchor=south] at (\s.north) {\iflangja{ステージ}{stage} \n};
  }
  \node[th] at (20,0) {};
  \foreach \y in {5,0,-5} { \node[th] at (60,\y) {}; }
  \node[th] at (100,0) {};
  \draw[->] (0,0) -- (s1);
  \foreach \i in {0,...,3} { \node[cell] at (35.5+3*\i,0) {}; }
  \foreach \i in {0,...,3} { \node[cell] at (75.5+3*\i,0) {}; }
  \draw[->] (s1) -- (34,0);  \draw[->] (45,0) -- (s2);
  \draw[->] (s2) -- (74,0);  \draw[->] (85,0) -- (s3);
  \draw[->] (s3) -- (116,0);
  \node[lab, anchor=north] at (40,-3) {\iflangja{バッファ}{buffer}};
  \node[lab, anchor=north] at (80,-3) {\iflangja{バッファ}{buffer}};
  \node[lab, anchor=north] at (60,-9) {\iflangja{スレッドプール}{thread pool}};
\end{tikzpicture}
   (width ~116 mm)
\begin{tikzpicture}[x=1mm, y=1mm, font=\footnotesize\sffamily,
  box/.style={draw, rounded corners=2pt, fill=ln-box, align=center, inner sep=1pt},
  th/.style={draw, circle, fill=white, minimum size=4mm, inner sep=0pt},
  cell/.style={draw, fill=white, minimum width=3mm, minimum height=5mm, inner sep=0pt},
  >={Stealth[length=1.6mm]},
  lab/.style={font=\scriptsize\sffamily, text=ln-muted}]
  \node[box, minimum width=14mm, minimum height=17mm] (s1) at (20,0) {};
  \node[box, minimum width=14mm, minimum height=17mm] (s2) at (60,0) {};
  \node[box, minimum width=14mm, minimum height=17mm] (s3) at (100,0) {};
  \foreach \s/\n in {s1/1,s2/2,s3/3} {
    \node[font=\scriptsize\sffamily, anchor=south] at (\s.north) {\iflangja{ステージ}{stage} \n};
  }
  \node[th] at (20,0) {};
  \foreach \y in {5,0,-5} { \node[th] at (60,\y) {}; }
  \node[th] at (100,0) {};
  \draw[->] (0,0) -- (s1);
  \foreach \i in {0,...,3} { \node[cell] at (35.5+3*\i,0) {}; }
  \foreach \i in {0,...,3} { \node[cell] at (75.5+3*\i,0) {}; }
  \draw[->] (s1) -- (34,0);  \draw[->] (45,0) -- (s2);
  \draw[->] (s2) -- (74,0);  \draw[->] (85,0) -- (s3);
  \draw[->] (s3) -- (116,0);
  \node[lab, anchor=north] at (40,-3) {\iflangja{バッファ}{buffer}};
  \node[lab, anchor=north] at (80,-3) {\iflangja{バッファ}{buffer}};
  \node[lab, anchor=north] at (60,-9) {\iflangja{スレッドプール}{thread pool}};
\end{tikzpicture}
   (width ~116 mm)
\begin{tikzpicture}[x=1mm, y=1mm, font=\footnotesize\sffamily,
  box/.style={draw, rounded corners=2pt, fill=ln-box, align=center, inner sep=1pt},
  th/.style={draw, circle, fill=white, minimum size=4mm, inner sep=0pt},
  cell/.style={draw, fill=white, minimum width=3mm, minimum height=5mm, inner sep=0pt},
  >={Stealth[length=1.6mm]},
  lab/.style={font=\scriptsize\sffamily, text=ln-muted}]
  \node[box, minimum width=14mm, minimum height=17mm] (s1) at (20,0) {};
  \node[box, minimum width=14mm, minimum height=17mm] (s2) at (60,0) {};
  \node[box, minimum width=14mm, minimum height=17mm] (s3) at (100,0) {};
  \foreach \s/\n in {s1/1,s2/2,s3/3} {
    \node[font=\scriptsize\sffamily, anchor=south] at (\s.north) {\iflangja{ステージ}{stage} \n};
  }
  \node[th] at (20,0) {};
  \foreach \y in {5,0,-5} { \node[th] at (60,\y) {}; }
  \node[th] at (100,0) {};
  \draw[->] (0,0) -- (s1);
  \foreach \i in {0,...,3} { \node[cell] at (35.5+3*\i,0) {}; }
  \foreach \i in {0,...,3} { \node[cell] at (75.5+3*\i,0) {}; }
  \draw[->] (s1) -- (34,0);  \draw[->] (45,0) -- (s2);
  \draw[->] (s2) -- (74,0);  \draw[->] (85,0) -- (s3);
  \draw[->] (s3) -- (116,0);
  \node[lab, anchor=north] at (40,-3) {\iflangja{バッファ}{buffer}};
  \node[lab, anchor=north] at (80,-3) {\iflangja{バッファ}{buffer}};
  \node[lab, anchor=north] at (60,-9) {\iflangja{スレッドプール}{thread pool}};
\end{tikzpicture}

  \caption{A three-stage pipeline.  Stages exchange requests through
    buffers; the slow second stage is a pool of three threads.}
  \label{fig:l03-pipeline}
\end{figure}
The pipeline offers specialization and locality, since every thread
repeatedly performs the same subtask; its costs are the effort of keeping
the stages balanced and the synchronization overhead of the buffers.

\begin{example}
  A three-stage pipeline has stage times \SI{10}{\milli\second},
  \SI{40}{\milli\second} and \SI{20}{\milli\second}.  With one thread per
  stage, \cref{eq:l03-pipeline-throughput} gives
  $\min(100, 25, 50) = 25$ requests per second, limited by stage~2.  With
  four threads in stage~2 the throughput becomes $\min(100, 100, 50) = 50$
  requests per second, and stage~3 is now the bottleneck.
\end{example}

\subsection{Layered}

In the \term{layered pattern}[階層化パターン] related subtasks are grouped
into layers, each served by its own threads, and every request passes down
through all the layers and back up (\cref{fig:l03-layered}).  Like the
pipeline it provides
specialization, while being less fine-grained.  It is not suitable for all
applications, and it requires synchronization between each layer and the
layers above and below it.

\begin{figure}[htbp]
  \centering
  % Layered pattern: a request passes down through the layers and back up.
% Usage: % Layered pattern: a request passes down through the layers and back up.
% Usage: % Layered pattern: a request passes down through the layers and back up.
% Usage: \input{figures/l03-layered}   (width ~84 mm)
\begin{tikzpicture}[x=1mm, y=1mm, font=\footnotesize\sffamily,
  box/.style={draw, rounded corners=2pt, fill=ln-box, align=center, inner sep=1pt},
  th/.style={draw, circle, fill=white, minimum size=4mm, inner sep=0pt},
  >={Stealth[length=1.6mm]},
  lab/.style={font=\scriptsize\sffamily, text=ln-muted}]
  \foreach \y/\n in {0/1,-8/2,-16/3} {
    \node[box, minimum width=52mm, minimum height=6.5mm] at (42,\y) {};
    \node[font=\scriptsize\sffamily, anchor=west] at (18,\y) {\iflangja{層}{layer} \n};
    \foreach \x in {44,54,64} { \node[th] at (\x,\y) {}; }
  }
  \draw[->] (12,5) -- (12,-23);
  \draw[->] (12,-23) -- (72,-23) -- (72,5);
  \node[lab, anchor=east] at (11,-8) {\iflangja{要求}{request}};
  \node[lab, anchor=west] at (73,-8) {\iflangja{応答}{response}};
\end{tikzpicture}
   (width ~84 mm)
\begin{tikzpicture}[x=1mm, y=1mm, font=\footnotesize\sffamily,
  box/.style={draw, rounded corners=2pt, fill=ln-box, align=center, inner sep=1pt},
  th/.style={draw, circle, fill=white, minimum size=4mm, inner sep=0pt},
  >={Stealth[length=1.6mm]},
  lab/.style={font=\scriptsize\sffamily, text=ln-muted}]
  \foreach \y/\n in {0/1,-8/2,-16/3} {
    \node[box, minimum width=52mm, minimum height=6.5mm] at (42,\y) {};
    \node[font=\scriptsize\sffamily, anchor=west] at (18,\y) {\iflangja{層}{layer} \n};
    \foreach \x in {44,54,64} { \node[th] at (\x,\y) {}; }
  }
  \draw[->] (12,5) -- (12,-23);
  \draw[->] (12,-23) -- (72,-23) -- (72,5);
  \node[lab, anchor=east] at (11,-8) {\iflangja{要求}{request}};
  \node[lab, anchor=west] at (73,-8) {\iflangja{応答}{response}};
\end{tikzpicture}
   (width ~84 mm)
\begin{tikzpicture}[x=1mm, y=1mm, font=\footnotesize\sffamily,
  box/.style={draw, rounded corners=2pt, fill=ln-box, align=center, inner sep=1pt},
  th/.style={draw, circle, fill=white, minimum size=4mm, inner sep=0pt},
  >={Stealth[length=1.6mm]},
  lab/.style={font=\scriptsize\sffamily, text=ln-muted}]
  \foreach \y/\n in {0/1,-8/2,-16/3} {
    \node[box, minimum width=52mm, minimum height=6.5mm] at (42,\y) {};
    \node[font=\scriptsize\sffamily, anchor=west] at (18,\y) {\iflangja{層}{layer} \n};
    \foreach \x in {44,54,64} { \node[th] at (\x,\y) {}; }
  }
  \draw[->] (12,5) -- (12,-23);
  \draw[->] (12,-23) -- (72,-23) -- (72,5);
  \node[lab, anchor=east] at (11,-8) {\iflangja{要求}{request}};
  \node[lab, anchor=west] at (73,-8) {\iflangja{応答}{response}};
\end{tikzpicture}

  \caption{The layered pattern: each layer groups related subtasks and has
    its own threads; a request passes down through the layers and back
    up.}
  \label{fig:l03-layered}
\end{figure}

\subsection{Comparing boss-workers and pipeline}

Which pattern completes a batch of requests faster depends on the number of
requests.\margin{Lesson~6 measures these patterns in real web servers:
Apache's worker MPM is boss-workers with one pool per process
(\texttt{ThreadsPerChild} = 25 by default),
and the event-driven Flash server handles many requests in a single
thread.}  Consider $n$ requests, each needing the same total processing
time.  A boss-workers design with $w$ workers and a negligible boss time
$t_{\mathrm{boss}}$ completes requests in groups of $w$, each group taking
the full task time $t$; a pipeline of $k$ equal stages of time $s$ completes
the first request after $k$ stage times and then one request per stage time:
\begin{equation}
  T_{\mathrm{bw}}(n) = \Bigl\lceil \frac{n}{w} \Bigr\rceil\, t ,
  \qquad
  T_{\mathrm{pl}}(n) = (k + n - 1)\, s .
  \label{eq:l03-batch-time}
\end{equation}

\begin{example}
  Both designs use four threads, and a request needs
  \SI{100}{\milli\second} of work.\caution{The formulas compare the
  completion time of a whole batch.  A pipeline does not shorten a single
  request: it still passes through all $k$ stages, and it waits in a buffer
  whenever the next stage is busy.}  The boss-workers design has one boss and
  $w=3$ workers with $t=\SI{100}{\milli\second}$; the pipeline has $k=4$
  stages of $s=\SI{25}{\milli\second}$.  Boss-workers completes 3 requests
  per \SI{100}{\milli\second} and the full pipeline 4, so boss-workers can
  lead only in small batches, where the pipeline's start-up time $k\,s$
  still counts.  Even there the outcome depends on $n$, because
  $T_{\mathrm{bw}}$ of \cref{eq:l03-batch-time} grows in steps: boss-workers
  is faster for $n = 2$, 3 and 6, the pipeline for $n = 4$, 7 and 8, the
  two tie for $n = 1$, 5 and 9, and for every $n \geq 10$ the pipeline
  wins:
  \begin{center}
    \small
    \begin{tabular}{@{}rrr@{}}
      \toprule
      $n$ & Boss-workers (\si{\milli\second}) & Pipeline (\si{\milli\second}) \\
      \midrule
       3 & 100 & 150 \\
       4 & 200 & 175 \\
       6 & 200 & 225 \\
      12 & 400 & 375 \\
      \bottomrule
    \end{tabular}
  \end{center}
\end{example}

\begin{keypoint}[Lesson summary]
  A thread is an execution context that shares its process's address space;
  threads provide parallelism, specialization and the use of idle time
  (\cref{eq:l03-idle}) more cheaply than processes.  Shared data must be
  protected with mutexes, and threads wait for conditions with condition
  variables, always re-testing the condition in a \texttt{while} loop.
  Proxy variables build richer policies such as readers-writers; signalling
  outside the mutex avoids spurious wake-ups; a consistent lock order avoids
  deadlocks.  User-level threads run on kernel-level threads according to
  the one-to-one, many-to-one or many-to-many model, and applications are
  structured as boss-workers, pipelines or layers.
\end{keypoint}

\end{document}
